\section{Relationship to prior work}\label{sec:related}
\subsection{Logic-based outer approximation and perspective cuts}
The extended hull construction belongs to disjunctive programming \citep{balas1998}; its convex nonlinear counterparts use perspective functions to represent bounded unions of convex sets \citep{ceria1999,grossmann2003}. Hull formulations of separate disjunctions are then intersected with global constraints and logic. That intersection need not be the convex hull of the complete GDP feasible set. This distinction matters when describing the relaxation targeted by our separation routines.

Logic-based outer approximation already combines selected-structure NLP subproblems with a linearized disjunctive master \citep{turkay1996}. Omitting inactive nonlinear constraints from those subproblems, retaining logical structure, and applying a hull transform to affine approximations are inherited ingredients. Our radial policy chooses additional linearization points without solving an optimization problem at each point, but the complete tested algorithm still uses NLP subproblems.

\citet{frangioni2006} derive perspective cuts for convex on/off mixed-integer programs. For a differentiable convex row $g(x)\leq0$, a tangent at $z$ gives
\begin{equation}\label{eq:related-perspective}
 \nabla g(z)^T v+
 \bigl(g(z)-\nabla g(z)^Tz\bigr)\lambda\leq0,
\end{equation}
where $v$ is a disaggregated variable and $\lambda$ its selection weight. For positive $\lambda$, this is precisely a linearization of $\lambda g(v/\lambda)$. \citet{bestuzheva2023} establish the same equivalence through affine cut extension and study perspective cuts in a general-purpose solver, including extensions based on valid nonconvex estimators. Accordingly, constructing coefficients in original-variable space and attaching an indicator does not itself distinguish a new cut family. Our hull ECP and ESH variants choose different $z$ in \eqref{eq:related-perspective}; the big-$M$ variants deactivate tangents in the original variables using valid big-$M$ coefficients.

\subsection{Radial supporting planes and disjunctive strengthening}
The point-tangent comparison follows the ECP tradition for convex mixed-integer nonlinear programming \citep{westerlund1995}. \citet{veinott1967} already selects supporting points by intersecting the boundary with a segment from a fixed interior point to the current exterior relaxation solution. \citet{kronqvist2016} develop ESH for convex mixed-integer nonlinear programming, using interior points, boundary searches, and an initial LP phase before mixed-integer refinement. \citet{serrano2020} place this construction in the earlier supporting-hyperplane tradition and show its relationship to Kelley's method applied to a gauge representation. Their result explains why a change in constraint representation can change point tangents while preserving the underlying feasible-set geometry. Our representation diagnostic makes a restricted version of that mechanism explicit and executable. It is not a new geometric foundation for ESH.

The SHOT implementation of \citet{lundell2022} is a close algorithmic precedent: it combines ESH and ECP separation with fixed-integer NLP heuristics, supports both single-tree callbacks and repeated-master execution, and falls back to ECP when no strict interior point is available. Its experiments include comparisons of tree policies. Our study retains GDP term structure and isolates radial versus point separation under matched settings within one GDP solver. SHOT also appears among our reformulation-based baselines; that comparison is distinct from the matched policy experiment.

The closest disjunctive ESH precedent is \citet{kronqvist2021}. Their strengthening procedure holds a cut normal fixed and solves separate convex optimization problems to obtain coefficients for the terms of a disjunction. They integrate these inequalities with ESH and compare against hull formulations while avoiding direct nonlinear perspective evaluation. Our hull variant instead maintains disaggregated variables and constructs term-specific tangents by line search in the original bounded domains. The normals may differ among terms, and the same radial tangent policy is also used with our original-variable big-$M$ master. This is a difference in separation procedure and master representation, not a demonstrated dominance relation. In particular, combining ESH with disjunctions or avoiding nonlinear perspective division cannot serve as a novelty claim for this paper.

\citet{trespalacios2016} construct cuts from a hull formulation strengthened by basic steps and project the resulting information into a big-$M$ model. Their separation uses optimization over a strengthened relaxation. Basic steps can incorporate interactions among disjunctions that are absent from separate-disjunction hulls. Our individual-row separation has a different work requirement and targets a different relaxation; it does not replace or subsume that method. No empirical comparison with their strengthening procedure is reported here.

\subsection{Exact conic formulations and conic outer approximation}
A nonquadratic expression does not establish that an exact conic hull is unavailable. \citet{bernal2024} formulate convex GDP alternatives with cones and obtain perspective-free hull representations under their stated assumptions. Their scope includes nonlinearities beyond quadratics. For quadratic rows, the revised manuscript of \citet{gusev2026}, specifically version~2 dated 17 March 2026, presents the conic exact hull reformulation (CEHR). For a row $g(x)=x^TQx+c^Tx+d\leq0$ with $Q\succeq0$, a nonnegative auxiliary variable $t$ gives the lifted constraints
\begin{equation}\label{eq:related-cehr}
 v^TQv\leq t\lambda,\qquad
 t+c^Tv+d\lambda\leq0,\qquad t\geq0,
\end{equation}
together with scaled bounds and $0\leq\lambda\leq1$. This has the same closed bounded-domain perspective target as the complete family of tangent cuts for the corresponding quadratic row. Comparing a finite outer approximation against a weaker algebraic alternative would therefore be insufficient; our quadratic comparison uses the exact conic formulation.

\citet{coey2020} develop outer approximation with conic certificates, including numerical tolerance issues and both iterative and solver-driven implementations. Their study distinguishes separation-only variants from algorithms using continuous conic subproblems. Thus affine master refinement, optional continuous optimization oracles, and single-tree integration are established algorithmic choices. The cone references in our study check supported hull roots and small fixed-assignment problems, and the quadratic conic baseline supplies a direct computational alternative. They do not constitute a complete comparison against a general conic outer-approximation solver.

\subsection{Scope of the contribution}
The literature establishes the constituent cut families and algorithmic ideas. Our original empirical result is the measured radial-versus-point tradeoff under matched GDP policies, together with its dependence on formulation, tree execution, incumbent recovery, and model scope. The accompanying proofs and diagnostic constructions make the interpretation checkable without claiming a new general convergence theory, a stronger limiting hull, or a uniformly stronger separator. A supporting plane can cut a different region from a point tangent without dominating it.

Recent related work also limits broad claims of a new GDP solution paradigm. \citet{nguyen2026} extend logic-based and cutting-plane approaches to infinite-dimensional GDP, a different problem class. Current \texttt{discopt} documentation describes algebraic ESH with interior-point root searches and a separate GDP reformulation choice \citep{discopt2026}. This is relevant implementation context, rather than peer-reviewed priority or an executed baseline in our study. It motivates keeping the claim specific to the present comparison. We make no claim that this is the first possible combination of GDP structure and radial separation, and the search for related work does not prove the absence of an equivalent implementation.
